\documentclass[12pt]{article}
\usepackage[utf8]{inputenc}
\usepackage[T1]{fontenc}
\usepackage{hyperref}
\title{Western Sisaala Language Resources for K-12 Curriculum Development}
\author{Compiled by Claude Sonnet 5.5 (claude-sonnet-5-5) for Andr\'as Kornai}
\date{}
\begin{document}
\maketitle

\section{Introduction}
Western Sisaala (ISO 639-3 ssl), also called Lambishi Sisaala, is a Gur language of northern Ghana, spoken around Tumu in the Upper West Region and in the Gonja area. It is closely related to Pasaale (Paasaal) and Tumulung Sisaala, and to Sisaali in Burkina Faso. Sources disagree on whether these are separate languages or dialects. Much of the material found is for Sisaala generally or for Pasaale, and it cannot be assumed to cover Western Sisaala.

\section{Primary Educational Resources}
No primers or textbooks specific to Western Sisaala turned up in the search.

\section{Linguistic and Reference Materials}
\begin{itemize}
\item Wikipedia, Sisaala language. It mentions a Sisaala--English dictionary (GILLBT, 1975): \url{https://en.wikipedia.org/wiki/Sisaala_language}
\item Sisaala Dictionary on Webonary (GILLBT and SILDEP). The search did not say which variety it covers: \url{https://www.webonary.org/sisaala/help/}
\item Omniglot, Paasaal (Pasaale Sisaala). It describes a Latin-script writing system developed by GILLBT for Pasaale: \url{https://omniglot.com/writing/paasaal.htm}
\end{itemize}

\section{Religious and Community Resources}
\begin{itemize}
\item Joshua Project, Western Sisaala in Ghana: \url{https://joshuaproject.net/index.php/people_groups/14953}
\end{itemize}
The search did not confirm a Bible translation specifically for Western Sisaala.

\section{Digital and Audio Resources}
Apart from the Webonary dictionary above, no audio, video or software resources turned up in the search.

\section{Curriculum Development Context}
GILLBT has done literacy and orthography work in the wider Sisaala group, so it is a natural contact. It is not established that any of this work is usable for Western Sisaala as it stands. Adult literacy among the Pasaale is reported to be very low, below 2\%. The relationship between the varieties needs checking before any Pasaale or Sisaala materials are adapted.
\end{document}